
\section{Matrix completion}\label{sec:mc-l2}

A pressing challenge often encountered in data science applications is estimation and learning in the face of \emph{missing data}.
To elucidate how to tackle this challenge via spectral methods, we delve into the renowned matrix completion problem in this section,
followed by another application called tensor completion in Section~\ref{sec:tensor-completion}.


Imagine that one observes a small subset of the entries in a large unknown matrix
and seeks to fill in all missing entries. An archetypal example is collaborative filtering, where one aims to predict the users' preferences on a collection of products based on partially revealed user-product ratings.   See Figure~\ref{fig:completion} for an illustration.  The problem, often referred to as matrix completion, is apparently ill-posed in general, as there are (much) fewer measurements than the unknowns.

%This is indeed a factor model in Section~\ref{sec:PCA} with missing data.  

Fortunately, if the matrix of interest exhibits certain low-dimensional structure, then reliable recovery becomes feasible.
A commonly encountered example of this kind concerns the case when the target matrix enjoys a low-rank structure. Again, take collaborative filtering for example:  the user-product rating matrix  might be well explained by a relatively small number of latent factors connecting users' preferences with products' attributes, thus resulting in an approximately low-rank matrix.  Motivated by its fundamental importance,  recent years have witnessed a flurry of research activity in studying low-rank matrix completion \citep{candes2009exact, keshavan2010matrix, gross2011recovering};
see \citet{chen2018harnessing,davenport2016overview} for overviews of recent developments.
In the sequel, we present a simple yet effective approach enabled by the spectral method, originally proposed in~\citet{achlioptas2007fast,keshavan2010matrix}.

\begin{figure}
\begin{center}
\includegraphics[width=0.35\textwidth]{figures/matrix_completion}
\end{center}
\caption{Illustration of matrix completion, where each ``$\checkmark$'' stands for an observed entry and each ``$?$'' represents a missing entry.} \label{fig:completion}
\end{figure}


\subsection{Problem formulation and assumptions}
\label{sec:problem-formulation-MC}

Suppose that we are interested in estimating an $n_{1}\times n_{2}$  rank-$r$ matrix  $\bm{M}^{\star} = [{M}^{\star}_{i,j}]_{1\leq i\leq n_1, 1\leq j\leq n_2}$.
Without loss of generality, we assume
$$n_{1}\leq n_{2}.$$  Denote by  $\bm{M}^{\star}=\bm{U}^{\star}\bm{\Sigma}^{\star}\bm{V}^{\star\top}$ the SVD of $\bm{M}^{\star}$, where the columns of $\bm{U}^{\star}\in \mathbb{R}^{n_1\times r}$ (resp.~$\bm{V}^{\star}\in \mathbb{R}^{n_2\times r}$) are the left (resp.~right) singular vectors of $\bm{M}^{\star}$, and $\bm{\Sigma}^{\star}$ is a diagonal matrix whose diagonal entries are the singular values of $\bm{M}^{\star}$. We define the condition number of the matrix $\bm{M}^{\star}$ to be $\kappa \coloneqq \sigma_1(\bM^\star)/\sigma_r(\bM^\star)$.


To capture the presence of missing data, we introduce an index subset $\Omega\subseteq [n_1]\times [n_2]$,
such that each entry ${M}^{\star}_{i,j}$ is observed if and only if $(i,j)\in \Omega$.
The goal is to reconstruct the singular subspaces $\bm{U}^{\star}$ and $\bm{V}^{\star}$, as well as the  full matrix $\bm{M}^{\star}$, based on entries observed over the sampling set $\Omega$.



\paragraph{Random sampling.}
Apparently, not all sampling patterns admit reliable estimation. For instance, if $\Omega$ contains only entries in the top half of the matrix,
then there is in general no hope to predict the bottom half of the matrix. In order to allow for meaningful matrix completion,
this monograph focuses on a natural random observation model commonly adopted in the literature, as formulated below.
%
\begin{assumption}[Random sampling]
\label{assump:mc-bernoulli}
Each entry of $\bm{M}^{\star}$ is observed independently with probability $0<p<1$, namely,
each $(i,j)\in [ n_{1}] \times [n_{2}]$ is included in $\Omega$ independently with probability $p$.
\end{assumption}
%
Under this model, we shall view the expected number of observed entries---namely, $pn_1n_2$---as the sample size. In truth, as long as $p$ is not overly small,  the number of observed entries is expected to concentrate  around its mean $pn_1n_2$.



\paragraph{Incoherence conditions.}
Caution needs to be exercised, however,  that the random sampling model alone does not guarantee effective recovery of an arbitrary low-rank matrix $\bm{M}^{\star}$.
Consider, for example, the following rank-1 matrix $\bm{M}^{\star}$ containing a single nonzero entry:
%
\[
%{\small
\left[\begin{array}{cccc}
1 & 0 & \cdots & 0\\
0 & 0 & \cdots & 0\\
\vdots & \vdots & \ddots & \vdots\\
0 & 0 & \cdots & 0
\end{array}\right] .
%}
\]
%
If $p=o(1)$, then with probability $1-p=1-o(1)$, the sampling pattern will fail to include the nonzero entry $M_{1,1}^{\star}$,
thus ruling out the possibility of faithful matrix recovery.
Consequently,  one needs to make sure that the sampling pattern does not suppress too much useful information.
Towards this end, the pioneering work \citet{candes2009exact, candes2010NearOptimalMC} singled out an incoherence parameter that plays a vital role.
%
\begin{definition}
\label{assump:mc-incoherence}
The incoherence parameter $\mu$ of the matrix $\bm{M}^{\star}\in\mathbb{R}^{n_{1}\times n_{2}}$ is defined as
%
\begin{align*}
	\mu\coloneqq \max\left\{  \frac{ n_1 \|\bm{U}^{\star}\|_{2,\infty}^2 }{r} ,  \frac{ n_2 \|\bm{V}^{\star}\|_{2,\infty}^2 }{r} \right\} .
	%\\
	%\|\bm{V}^{\star}\|_{2,\infty} & \leq\sqrt{\mu/n_{2}}\,\|\bm{V}^{\star}\|_{\mathrm{F}}=\sqrt{\mu r/n_{2}}.
\end{align*}
%
\end{definition}
%
\begin{remark}
\label{remark:range-mu-MC}
%
Recognizing the following basic relation
%
\[
	\frac{r}{n_1}=\frac{1}{n_1} \|\bm{U}^{\star}\|_{\mathrm{F}}^2 \leq \|\bm{U}^{\star}\|_{2,\infty}^2 \leq \|\bm{U}^{\star}\|^2 = 1
\]
%
and an analogous one for $\bm{V}^{\star}$, we have $1\leq \mu \leq \max\{n_1,n_2\}/r= n_2/r$.
%
\end{remark}


In words, a small $\mu$ indicates that the energy of the singular vectors is spread out across different elements, namely,  the singular subspace of $\bm{M}^{\star}$ is not too ``aligned'' with any of the standard basis vectors,
thus ensuring that entrywise observations provide somewhat equalized information about the full spectrum of $\bm{M}^{\star}$.
The following lemma summarizes a few immediate consequences of this definition, with the proof deferred to Section~\ref{sec:proof-auxiliary-lemmas-MC-L2}.
%
\begin{lemma}\label{claim:mu-incoherence}
Assume that  $\bm{M}^{\star}\in\mathbb{R}^{n_{1}\times n_{2}}$ is $\mu$-incoherent. Then the following relations hold
%
\begin{subequations}
\label{subeq:mu-incoherence}
\begin{align}
	\|\bm{M}^{\star}\|_{2,\infty} & \leq\sqrt{\mu r/n_{1}}\, \big\| \bm{M}^{\star} \big\|; \quad
	\|\bm{M}^{\star\top}\|_{2,\infty}  \leq\sqrt{\mu r/n_{2}}\, \big\| \bm{M}^{\star} \big\|;\label{eq:mc-row-norm-upper}\\
	& \qquad\quad \|\bm{M}^{\star}\|_{\infty}  \leq\mu r  \big\| \bm{M}^{\star} \big\| /\sqrt{n_{1}n_{2}} . \label{eq:mc-entry-upper}
\end{align}
\end{subequations}
%
\end{lemma}



\paragraph{Additional notation.}

We find it convenient to introduce a Euclidean projection operator $\mathcal{P}_{\Omega}:\mathbb{R}^{n_{1}\times n_{2}}\mapsto\mathbb{R}^{n_{1}\times n_{2}}$ such that
%
\begin{equation}\label{defn:Pomega}
\big[\mathcal{P}_{\Omega}(\bm{A}) \big]_{i,j}=\begin{cases}
A_{i,j},\quad & \text{if }(i,j)\in\Omega \\
0, & \text{else}
\end{cases}
\end{equation}
%
for any matrix $\bm{A} =[A_{i,j}]\in\mathbb{R}^{n_{1}\times n_{2}}$.
With this notation in place, matrix completion amounts to recovering $\bm{M}^{\star}$ on the basis of $\mathcal{P}_{\Omega}(\bm{M}^{\star})$.




\subsection{Algorithm}\label{sec:mc-alg}

To apply the spectral method, the first step is to form a reasonable approximation $\bm{M}$ of the unknown matrix $\bm{M}^{\star}$.
By virtue of the random sampling model (cf.~Assumption~\ref{assump:mc-bernoulli}), a candidate approximation can be obtained from the observed data matrix via inverse probability weighting:
%
\begin{align}
	\bm{M}   \coloneqq p^{-1}\mathcal{P}_{\Omega}(\bm{M}^{\star}).
	\label{eq:rescaled-M-matrix-completion-l2}
\end{align}
%
The rationale is that $\bm{M}$ forms an unbiased estimate of the ground truth, namely, $$\mathbb{E}[\bm{M}]=\bm{M}^{\star},$$ where
the expectation is taken over the randomness in $\Omega$.


As a result, the proposed spectral method proceeds by computing the rank-$r$ SVD $\bm{U}\bm{\Sigma}\bm{V}^{\top}$ of the matrix $\bm{M}$ constructed in \eqref{eq:rescaled-M-matrix-completion-l2}, and employing $\bm{U}\in \mathbb{R}^{n_1\times r}$, $\bm{V}\in \mathbb{R}^{n_2\times r}$  and $\bm{U}\bm{\Sigma}\bm{V}^{\top}$ as estimates of $\bm{U}^{\star}$, $\bm{V}^{\star}$ and $\bm{M}^{\star}$, respectively.






\subsection{Performance guarantees}\label{sec:mc-l2-theory}

As before, whether the subspace $\bm{U}$ (resp.~$\bm{V}$) is close
to $\bm{U}^{\star}$ (resp.~$\bm{V}^{\star}$) relies crucially on
the size of the perturbation $\|\bm{M}-\bm{M}^{\star}\|$.
Therefore, we begin by developing an upper bound on this quantity; the proof is based on the matrix Bernstein inequality and is postponed to Section~\ref{sec:proof-auxiliary-lemmas-MC-L2}.

\begin{lemma}
\label{lemma:mc-noise-bound}
Consider the settings in Section~\ref{sec:problem-formulation-MC}. Suppose that $n_{2}p\geq C\mu r\log n_{2}$ for some
 constant $C>0$. Then with probability at least
	$1-O(n_{2}^{-10})$, the matrix $\bm{M}$ constructed in \eqref{eq:rescaled-M-matrix-completion-l2} obeys
%
\[
	\big\| \bm{M} -\bm{M}^{\star} \big\|
	\lesssim\sqrt{\frac{\mu r\log n_{2}}{n_{1}p}}\, \big\| \bm{M}^{\star} \big\|.
\]
%
\end{lemma}
%
With this perturbation bound in place, we are equipped to apply Wedin's $\sin\bm{\Theta}$ theorem to obtain the following results.  The condition on the sample size in Theorem~\ref{thm:mc-l2-subspace} is stronger than that in Lemma~\ref{lemma:mc-noise-bound}, as we need to control the eigengap in the following theorem.
%
\begin{theorem}
\label{thm:mc-l2-subspace}
Consider the settings in Section~\ref{sec:problem-formulation-MC}.
Suppose that $n_{1}p\geq C_{1}\kappa^{2}\mu r\log n_{2}$ for some sufficiently
large constant $C_{1}>0$. Then with probability exceeding $1-O(n_2^{-10})$,
%
\begin{align*}
\max\Big\{ \mathsf{dist}\left(\bm{U},\bm{U}^{\star}\right),\mathsf{dist}\left(\bm{V},\bm{V}^{\star}\right)\Big\}  & \lesssim\kappa\sqrt{\frac{\mu r \log n_{2}}{n_{1}p}} .
%;\\
%\max\Big\{ \mathsf{dist}_{\mathrm{F}}\left(\bm{U},\bm{U}^{\star}\right),\mathsf{dist}_{\mathrm{F}}\left(\bm{V},\bm{V}^{\star}\right)\Big\}  & \lesssim\kappa\sqrt{\frac{\mu r^{2} \log n_{2}}{n_{1}p}}.
\end{align*}
%
\end{theorem}
%
\begin{proof}
%	
As a direct consequence of Lemma \ref{lemma:mc-noise-bound}, one has
%
\[
	\big\| \bm{M} -\bm{M}^{\star} \big\|
	\lesssim\sqrt{\frac{\mu r\log n_{2}}{n_{1}p}} \,\big\| \bm{M}^{\star}\big\| \le \Big( 1 - \frac{1}{\sqrt{2}} \Big) \sigma_{r}(\bm{M}^{\star}),
\]
%
provided that $n_{1}p\geq C_{1}\kappa^{2}\mu r\log n_{2}$ for some  large enough constant $C_{1}>0$.
Apply Wedin's theorem (cf.~\eqref{eq:wedin-simpler-friendly}) and Lemma \ref{lemma:mc-noise-bound} to obtain
%
\begin{align*}
 & \max\Big\{ \mathsf{dist}\left(\bm{U},\bm{U}^{\star}\right),\mathsf{dist}\left(\bm{V},\bm{V}^{\star}\right) \Big\}
	\leq  \frac{2\big\|\bm{M} -\bm{M}^{\star}\big\|}{\sigma_{r}(\bm{M}^{\star})}
	%\\ % & \qquad\qquad
	%\lesssim\frac{\sqrt{\frac{\mu r\log n_{2}}{n_{1}p}}\,\big\|\bm{M}^{\star}\big\|}{\sigma_{r}(\bm{M}^{\star})}\asymp
	\lesssim \kappa\sqrt{\frac{\mu r\log n_{2}}{n_{1}p}}
\end{align*}
%
as claimed.
%The proof for the $\mathsf{dist}_{\mathrm{F}}(\cdot,\cdot)$ bounds follows from the same argument and is hence omitted.
%
\end{proof}


As an important implication of Theorem~\ref{thm:mc-l2-subspace}, once the sample size exceeds
$$pn_1n_2 \gg \kappa^{2}\mu r n_2  \log n_{2}, $$
then the spectral estimate achieves consistent estimation in the sense that
%
\[
	\max\Big\{ \mathsf{dist}\left(\bm{U},\bm{U}^{\star}\right),\mathsf{dist}\left(\bm{V},\bm{V}^{\star}\right)\Big\} = o(1).
\]
%
Given that $pn_1n_2\gtrsim \mu n_2 r \log n_2$ is an information-theoretic sampling requirement for reliable matrix completion when $r=o(n_1/\log n_2)$ \citep{candes2010NearOptimalMC},
Theorem~\ref{thm:mc-l2-subspace} confirms the near optimality of spectral methods---in terms of the scaling with $n_1$, $n_2$ and $p$---when it comes to consistent subspace estimation.



Before moving forward to the proof, we further characterize the statistical accuracy of $\bm{U}\bm{\Sigma}\bm{V}^{\top}$ in estimating the unknown matrix $\bm{M}^{\star}$. Accomplishing this only requires Lemma~\ref{lemma:mc-noise-bound}, without any need of the singular subspace perturbation theory. This result will also come in handy when we turn to discussing entrywise estimation accuracy in Chapter~\ref{cha:Linf-theory}.
%
\begin{theorem}
\label{thm:mc-l2-matrix}
Consider the settings in Section~\ref{sec:problem-formulation-MC}.
Suppose that $n_{2}p\geq C\mu r\log n_{2}$ for some sufficiently large constant
$C>0$. Then with probability at least $1-O(n_{2}^{-10})$, one has
%
\begin{align*}
%\|\bm{U}\bm{\Sigma}\bm{V}^{\top}-\bm{M}^{\star}\| & \lesssim\sqrt{\frac{\mu r}{n_{1}p}\log n_{2}}\,\sigma_{1}(\bm{M}^{\star});\\
\|\bm{U}\bm{\Sigma}\bm{V}^{\top}-\bm{M}^{\star}\|_{\mathrm{F}} & \lesssim\sqrt{\frac{\mu r^{2}\log n_{2}}{n_{1}p}}\, \big\| \bm{M}^{\star} \big\|.
\end{align*}
%
\end{theorem}
%
\begin{proof}
First, note
%
\begin{align*}
\|\bm{U}\bm{\Sigma}\bm{V}^{\top}-\bm{M}^{\star}\| &\leq\|\bm{U}\bm{\Sigma}\bm{V}^{\top}- \bm{M} \|+\|\bm{M} -\bm{M}^{\star}\| \leq2\|\bm{M}-\bm{M}^{\star}\|,
\end{align*}
%
where the first inequality comes from the triangle inequality, and the second inequality follows from the fact that $\bm{U}\bm{\Sigma}\bm{V}^{\top}$ is the best rank-$r$ approximation to $\bm{M}$,
i.e.,
$$\|\bm{U}\bm{\Sigma}\bm{V}^{\top}-\bm{M}\|=\min_{\bm{Z}: \mathsf{rank}(\bm{Z})\leq r}\|\bm{Z}-\bm{M}\| \leq \|\bm{M}-\bm{M}^{\star}\|. $$
Additionally, it is observed that $\bm{U}\bm{\Sigma}\bm{V}^{\top}-\bm{M}^{\star}$ has rank at most $2r$, which implies
%
\[
	\|\bm{U}\bm{\Sigma}\bm{V}^{\top}-\bm{M}^{\star}\|_{\mathrm{F}}\leq\sqrt{2r}\|\bm{U}\bm{\Sigma}\bm{V}^{\top}-\bm{M}^{\star}\| \leq 2\sqrt{2r} \big\|\bm{M} - \bm{M}^{\star} \big\|.
\]
%
This combined with Lemma~\ref{lemma:mc-noise-bound} immediately concludes the proof. \end{proof}

 


\subsection{Proof of auxiliary lemmas}
\label{sec:proof-auxiliary-lemmas-MC-L2}

\paragraph{Proof of Lemma~\ref{claim:mu-incoherence}.}
First of all, the  $\|\cdot\|_{2,\infty}$ norm of $\bm{M}^{\star}$
can be upper bounded by
%
\begin{equation*}
\|\bm{M}^{\star}\|_{2,\infty}=\|\bm{U}^{\star}\bm{\Sigma}^{\star}\bm{V}^{\star\top}\|_{2,\infty}\leq\|\bm{U}^{\star}\|_{2,\infty}\|\bm{\Sigma}^{\star}\| \, \|\bm{V}^{\star}\|\leq\sqrt{\frac{\mu r}{n_{1}}} \big\| \bm{M}^{\star} \big\|.
\end{equation*}
%
Here, the first inequality arises from the elementary bounds
$\|\bm{A}\bm{B}\|_{2,\infty}\leq\|\bm{A}\|_{2,\infty}\|\bm{B}\|$
and $\|\bm{A}\bm{B}\|\leq\|\bm{A}\|\, \|\bm{B}\|$, whereas the last relation
uses Definition~\ref{assump:mc-incoherence}, the orthonormality
of $\bm{V}^{\star}$, and identifies $\|\bm{\Sigma}^{\star}\|$ with
$\| \bm{M}^{\star} \|$. The  bound on $\|\bm{M}^{\star\top}\|_{2,\infty}$ can be derived analogously and is omitted for brevity.



In addition, the matrix $\bm{M}^{\star}$ is elementwise bounded by
%
\begin{equation*}
\|\bm{M}^{\star}\|_{\infty}=\|\bm{U}^{\star}\bm{\Sigma}^{\star}\bm{V}^{\star\top}\|_{\infty}\leq\|\bm{U}^{\star}\|_{2,\infty}\|\bm{V}^{\star}\|_{2,\infty} \|\bm{\Sigma}^{\star}\|
	\leq\frac{\mu r}{\sqrt{n_{1}n_{2}}} \big\| \bm{M}^{\star} \big\|.
\end{equation*}
%
Here, the first inequality follows from the fact $\|\bm{A}\bm{B}^{\top}\|_{\infty}\leq\|\bm{A}\|_{2,\infty}\|\bm{B}\|_{2,\infty}$
and the aforementioned one $\|\bm{A}\bm{B}\|_{2,\infty}\leq\|\bm{A}\|_{2,\infty}\|\bm{B}\|$,
while the last inequality again relies on Definition~\ref{assump:mc-incoherence}.
%\end{proof}


\paragraph{Proof of Lemma~\ref{lemma:mc-noise-bound}.} Note that the matrix $\bm{E}\coloneqq p^{-1}\mathcal{P}_{\Omega}(\bm{M}^{\star})-\bm{M}^{\star}$
can be expressed as the sum of $n_{1}n_{2}$ i.i.d.~random matrices
%
\[
  \frac{1}{p}  \mathcal{P}_{\Omega}(\bm{M}^{\star})-\bm{M}^{\star}=\sum_{i=1}^{n_{1}}\sum_{j=1}^{n_{2}}\underbrace{\big(p^{-1} \delta_{i,j} -1 \big)M_{i,j}^{\star}\bm{e}_{i}\bm{e}_{j}^{\top}}_{\eqqcolon\bm{X}_{i,j}}.
\]
%
Here, $\delta_{i,j}$ (which indicates whether the $(i,j)$-th entry is observed) follows an independent Bernoulli distribution
with parameter~$p$, and $\bm{e}_{i}$ stands for the $i$-th standard
basis vector of appropriate dimensions. It is easily seen that for each $(i,j)$,
%
\[
	\mathbb{E}[\bm{X}_{i,j}]=\bm{0}\quad\text{and}\quad\|\bm{X}_{i,j}\|\leq  \frac{1}{p} \|\bm{M}^{\star}\|_{\infty}\leq\frac{\mu r}{p\sqrt{n_{1}n_{2}}} \big\| \bm{M}^{\star} \big\|,
\]
%
where the last relation results from the entrywise upper bound (\ref{eq:mc-entry-upper})
on $\bm{M}^{\star}$. In order to apply the matrix Bernstein inequality (cf.~Corollary~\ref{thm:matrix-Bernstein-friendly}),
we need to control the variance statistic
%
\[
	v\coloneqq\max \Big\{ \Big\Vert \sum\nolimits _{i,j}\mathbb{E}\left[\bm{X}_{i,j}\bm{X}_{i,j}^{\top}\right]\Big\Vert ,\Big\Vert \sum\nolimits _{i,j}\mathbb{E}\left[\bm{X}_{i,j}^{\top}\bm{X}_{i,j}\right] \Big\Vert \Big\} .
\]
%
Regarding the first variance term, we have
%
\begin{align*}
\sum\nolimits _{i,j}\mathbb{E}\left[\bm{X}_{i,j}\bm{X}_{i,j}^{\top}\right]
	& =\sum\nolimits _{i,j}\mathbb{E}\left[ \big( p^{-1} \delta_{i,j} -1 \big)^{2}(M_{i,j}^{\star})^{2}\bm{e}_{i}\bm{e}_{j}^{\top}\bm{e}_{j}\bm{e}_{i}^{\top}\right]\\
 & =\frac{1-p}{p}\sum\nolimits _{i,j}\big(M_{i,j}^{\star}\big)^{2}\bm{e}_{i}\bm{e}_{i}^{\top}=\frac{1-p}{p}\sum_{i=1}^{n_{1}}\|\bm{M}_{i,\cdot}^{\star}\|_{2}^{2}\bm{e}_{i}\bm{e}_{i}^{\top} \\
	& \preceq \frac{1-p}{p}  \|\bm{M}^{\star}\|_{2,\infty}^{2} \bm{I}_{n_1} \preceq \frac{\mu r}{n_1p}  \|\bm{M}^{\star}\|^{2} \,\bm{I}_{n_1} .
\end{align*}
%
Here, the first identity arises from the definition of $\bm{X}_{i,j}$,
the second one calculates the variance of Bernoulli random variables,
and the last line relies on the upper bound~(\ref{eq:mc-row-norm-upper}).
Similarly, the second term in the variance statistic enjoys the following characterization:
%
\[
	\sum\nolimits _{i,j}\mathbb{E}\left[\bm{X}_{i,j}^{\top}\bm{X}_{i,j}\right] \preceq \frac{\mu r}{n_2p} \|\bm{M}^{\star}\|^{2} \,\bm{I}_{n_2}.
\]
%
Taking the above  relations together and recalling that $n_1\leq n_2$ give
%
\begin{align*}
v & %\leq\frac{1}{p}\max\left\{ \|\bm{M}^{\star}\|_{2,\infty}^{2},\|\bm{M}^{\star\top}\|_{2,\infty}^{2}\right\}
	% \leq\frac{1}{p}\left\{ \frac{\mu r}{n_{1}} \big\| \bm{M}^{\star} \big\|^2,  \frac{\mu r}{n_{2}} \big\| \bm{M}^{\star} \big\|^2 \right\} \\
	\leq\frac{\mu r}{n_{1}p} \big\| \bm{M}^{\star} \big\|^2.
\end{align*}
% 


With the above bounds in place, invoking matrix Bernstein (see Corollary~\ref{thm:matrix-Bernstein-friendly}) reveals that: with probability at least $1-O(n_2^{-10})$,
%
\begin{align*}
	\| \bm{E}\| & \lesssim \sqrt{\frac{\mu r \| \bm{M}^{\star} \|^2 \log n_{2}}{n_{1}p} }+\frac{\mu r \| \bm{M}^{\star} \| \log n_{2}}{p\sqrt{n_{1}n_{2}}}
	\asymp \sqrt{\frac{\mu r \| \bm{M}^{\star} \|^2 \log n_{2}}{n_{1}p} } ,
\end{align*}
%
where the last inequality is valid as long as  $n_{2}p \gtrsim \mu r\log n_{2}$.



